\subsection{Class II(E)}\label{sec:proof-s3}
\begin{proof}[Proof of \cref{prop:class-ii-e}]
We want plane-section traces \(\TraceDeDpRoot\) that also come from some
\(\mathbf{X}\) in the equilibrium-invariance group $G$. 
To this end, fix a reference plane-section trace \(\TraceDeDpRoot^R\) of the form \eqref{eq:T0-plane-structure}, with blocks
\(\TraceBT^R, \TraceFT^R,\dots\), and look at
\begin{equation*}
  \TraceDeDpRoot = \mathbf{X}\,\TraceDeDpRoot^R.
\end{equation*}

Equilibrium invariance requires:
\begin{equation*}
  \TraceDeDpRoot^\top \mathbb{J}\,\TraceDeDpRoot^{-\top}
  = \TraceDeDpRoot^{R,\top} \mathbb{J}\,\TraceDeDpRoot^{R,-\top}
  \;\Longleftrightarrow\;
  \mathbf{X}^\top \mathbb{J}\,\mathbf{X}^{-\top} = \mathbb{J}.
\end{equation*}
So \(\mathbf{X}^\top \in G\).
Write \(\mathbf{X}^\top = \mathbf{K}\) with the structure \eqref{eq:K-structure}, so
\begin{equation}
  \mathbf{X} = \mathbf{K}^\top =
  \begin{bmatrix}
    K_{11}     & \mathbf{0}^\top      & K_{31}      & \bm{k}_{41}^\top \\
    \bm{k}_{12}& \mathbf{K}_{22}^\top & \bm{k}_{32} & \mathbf{K}_{42}^\top \\
    K_{13}     & \mathbf{0}^\top      & K_{33}      & \bm{k}_{43}^\top \\
    \mathbf{0} & \mathbf{0}           & \mathbf{0}  & \mathbf{K}_{44}^\top
  \end{bmatrix}.
  \tag{4}
\end{equation}

Now impose: \(\TraceDeDpRoot = \mathbf{X}\,\TraceDeDpRoot^R\) must again have the
plane-section pattern \eqref{eq:T0-plane-structure}. 
Multiply in blocks and enforce only the fixed
zeros/ones of \eqref{eq:T0-plane-structure}.

Let \(\TraceBT^R\) be the \(2\times 1\) block in \(\TraceDeDpRoot^R\). 
Computing the blocks of \(\TraceDeDpRoot\) and enforcing:

\begin{itemize}
  \item \(T_{11} = 1\) gives \(K_{11} = 1\).
  \item \(T_{21} = \mathbf{0}\) gives \(\bm{k}_{12} = \mathbf{0}\).
  \item \(T_{22} = \mathbf{0}\) gives \(\mathbf{K}_{42} = \mathbf{0}\).
  \item \(T_{31} = 0\) gives \(K_{13} = 0\).
  \item \(T_{32} = \mathbf{0}\) gives \(\bm{k}_{43} = \mathbf{0}\).
  \item \(T_{33} = 1\) gives \(K_{33} = 1\).
  \item \(T_{42} = -\FrameBasis^\times\) yields
    \[
      \mathbf{K}_{44}^\top (-\FrameBasis^\times) = -\FrameBasis^\times
      \;\Rightarrow\;
      \mathbf{K}_{44} = \oneS,
    \]
    which also forces \(\mathbf{K}_{22} = \oneS\).
  \item \(T_{13} = 0\) gives the single scalar constraint
    \begin{equation*}
      K_{31} + \bm{k}_{41}\cdot\TraceBT^R = 0.
      % \tag{5}
    \end{equation*}
\end{itemize}

All other blocks of \(\TraceDeDpRoot\) are allowed to be arbitrary, and do not impose further equations. 
Collecting the conditions on \(\mathbf{K}\):
\[
  K_{11} = 1,\quad
  K_{13} = 0,\quad
  K_{33} = 1,\quad
  \bm{k}_{12} = \mathbf{0},\quad
  \mathbf{K}_{22} = \oneS,\quad
  \mathbf{K}_{42} = \mathbf{0},\quad
  \bm{k}_{43} = \mathbf{0},\quad
  \mathbf{K}_{44} = \oneS,
\]
with
\[
  K_{31} = -\bm{k}_{41}\cdot\TraceBT^R.
\]

The only free blocks left in \(\mathbf{K}\) are the two \(2\)-vectors
\(\bm{k}_{32},\bm{k}_{41}\).  So an admissible plane-section /
equilibrium-invariant \(\mathbf{K}\in G\) has block form
\begin{equation*}
  \mathbf{K} =
  \begin{bmatrix}
    1 & \mathbf{0}^\top       & 0 & \mathbf{0}^\top \\
    \mathbf{0} & \oneS        & \mathbf{0} & \mathbf{0} \\
    -\bm{k}_{41}\cdot\TraceBT^R & \bm{k}_{32}^\top & 1 & \mathbf{0}^\top \\
    \bm{k}_{41} & \mathbf{0}  & \mathbf{0} & \oneS
  \end{bmatrix}
\end{equation*}
with \(\bm{k}_{32},\bm{k}_{41}\in\mathbb{R}^2\) arbitrary. 
The corresponding admissible plane-section \(\mathbf{X}\) (since
\(\mathbf{X}^\top = \mathbf{K}\)) is
\begin{equation*}
  \mathbf{X} =
  \begin{bmatrix}
    1 & \mathbf{0}^\top & -\bm{k}_{41}\cdot\TraceBT^R & \bm{k}_{41}^\top \\
    \mathbf{0} & \oneS   & \bm{k}_{32} & \mathbf{0} \\
    0 & \mathbf{0}^\top  & 1 & \mathbf{0}^\top \\
    \mathbf{0} & \mathbf{0} & \mathbf{0} & \oneS
  \end{bmatrix}
\end{equation*}
Thus, P1 and S2 picks out exactly this 4‑parameter
subgroup inside \(G\). 
%
It is convenient to rename
\[
  \bm{b} \coloneq \bm{k}_{41},\qquad
  \bm{c} \coloneq \bm{k}_{32}.
\]
Then
\begin{equation*}
  \mathbf{X} =
  \begin{bmatrix}
    1 & \mathbf{0} & -\bm{b}\cdot {\TraceBT}^R & \bm{b}^\top \\
    \mathbf{0} & \oneS  & \bm{c} & \mathbf{0} \\
    0 & \mathbf{0} & 1 & \mathbf{0} \\
    \mathbf{0} & \mathbf{0} & \mathbf{0} & \oneS
  \end{bmatrix},
\end{equation*}
%
Write the reference plane-section trace as
\begin{equation*}
  \TraceDeDpRoot^R =
  \begin{pmatrix}
    1          & \bm{r}_R^\top & 0 & {\TraceNF}_R^\top \\
    \mathbf{0} & \mathbf{0}    & {\TraceFT}_R & {\TraceFF}_R \\
    0          & \mathbf{0} & 1 & {\TraceTF}_R^\top \\
    \mathbf{0} & -\FrameBasis^\times & {\TraceBT}_R & {\TraceMF}_R
  \end{pmatrix}.
\end{equation*}
Then \(\TraceDeDpRoot = \mathbf{X}\,\TraceDeDpRoot^R\) has the same plane-section
pattern
\begin{equation*}
  \TraceDeDpRoot =
  \begin{bmatrix}
    1          & \CenterTrace^\top & 0 & \TraceNF^\top \\
    \mathbf{0} & \mathbf{0}              & \TraceFT & \TraceFF \\
    0          & \mathbf{0}              & 1 & \TraceTF^\top \\
    \mathbf{0} & -\FrameBasis^\times     & \TraceBT & \TraceMF
  \end{bmatrix},
\end{equation*}
with the new blocks expressed in terms of the reference ones and
\(\bm{b},\bm{c}\):

\begin{align*}
  \CenterTrace^\top
    &= \bm{r}_R^\top - \bm{b}^\top \FrameBasis^\times, \\ %[4pt]
  \TraceFT
    &= {\TraceFT}_R + \bm{c}, \\ %[4pt]
  \TraceFF
    &= {\TraceFF}_R + \bm{c}\otimes{\TraceTF}_R, \\ %[4pt]
  \TraceNF^\top
    &= {\TraceNF}_R^\top
       - (\bm{b}^\top {\TraceBT}_R)\,{\TraceTF}_R^\top
       + \bm{b}^\top {\TraceMF}_R, \\%[4pt]
  \TraceTF
    &= {\TraceTF}_R, \\%[4pt]
  \TraceBT
    &= {\TraceBT}_R, \\%[4pt]
  \TraceMF
    &= {\TraceMF}_R.
\end{align*}
%
Define
\[
  \Delta \TraceDeDpRoot \coloneq \TraceDeDpRoot - \TraceDeDpRoot^R =
  \begin{bmatrix}
    0 & \Delta \bm{r}^\top & 0 & \Delta \TraceNF^\top \\[2pt]
    \mathbf{0} & \mathbf{0} & \Delta \TraceFT & \Delta \TraceFF \\[2pt]
    0 & \mathbf{0}^\top & 0 & \mathbf{0}^\top \\[2pt]
    \mathbf{0} & \mathbf{0} & \mathbf{0} & \mathbf{0}
  \end{bmatrix},
\quad\text{with}\quad
\]
\begin{align*}
  \Delta \bm{r}
    &= \CenterTrace - \bm{r}_R
     = \FrameBasis\times\bm{b}, \\ %[4pt]
  \Delta \TraceNF^\top
    &= \TraceNF^\top - {\TraceNF}_R^\top
     = -({\TraceBT}_R \cdot \bm{b})\,{\TraceTF}_R^\top
       + \bm{b}^\top {\TraceMF}_R, \\ %[4pt]
  \Delta \TraceFT
    &= \TraceFT - {\TraceFT}_R = \bm{c}, \\ %[4pt]
  \Delta \TraceFF
    &= \TraceFF - {\TraceFF}_R = \bm{c}\otimes {\TraceTF}_R.
\end{align*}
The remaining blocks
\(\TraceTF,\TraceBT,\TraceMF\) are unchanged. 
In summary:

\begin{itemize}
  \item \(\bm{b}\) is a \emph{trace-center gauge} for axial–bending:
        it shifts the trace center \(\CenterTrace\) and induces a
        matching change in \(\TraceNF\) while preserving equilibrium and
        energy structure.
  \item \(\bm{c}\) is a \emph{shear-basis gauge}: it mixes shear with
        torsion/flexure via \(\TraceFT,\TraceFF\) without affecting the
        purely torsional/bending blocks.
  \item The blocks corresponding to pure torsion and pure bending
        (the 1’s in columns 1 and 3, the \(-\FrameBasis^\times\) block,
        and \(\TraceTF,\TraceBT,\TraceMF\)) are structural invariants of
        the equilibrium-invariant plane-section family.
\end{itemize}
\end{proof}